\documentclass{article}

% NeurIPS 2026 style, preprint mode (not a NeurIPS submission; footer overridden below).
\usepackage[preprint]{neurips_2026}

\usepackage[utf8]{inputenc}
\usepackage[T1]{fontenc}
\usepackage{hyperref}
\usepackage{url}
\usepackage{booktabs}
\usepackage{amsmath,amssymb,amsfonts}
\usepackage{nicefrac}
\usepackage{microtype}
\usepackage{xcolor}
\usepackage{graphicx}
\usepackage{tabularx}
\usepackage{multirow}
\usepackage{caption}
\usepackage{float}
\captionsetup{font=small,labelfont=bf}

\hypersetup{pdfauthor={Gunnar Andra-Thomas},pdftitle={Long-Horizon Predictive Subspaces in Mouse Secondary Motor Cortex Track Time in the Session, Not the Block Prior},colorlinks=true,linkcolor=blue!50!black,citecolor=blue!50!black,urlcolor=blue!50!black}

% Footer: this is a course project, not a NeurIPS paper.
\makeatletter
\renewcommand{\@noticestring}{CSE 8803 NBD final project, Fall 2026}
\makeatother

\graphicspath{{figures/}}

% Figure with a framed placeholder if the PDF is missing (lets the document compile before figures exist).
% report/figures/rfig_*.pdf are built by report/scripts/make_figures.py from the result files.
\newcommand{\maybefig}[3]{%
  \IfFileExists{figures/#1}{\includegraphics[#2]{#1}}{%
    \fbox{\parbox[c][#3][c]{0.9\linewidth}{\centering\small Figure file \texttt{#1} not yet available.}}}}

\newcommand{\RR}{\mathbb{R}}
\newcommand{\tr}{\operatorname{tr}}
\newcommand{\orth}{\operatorname{orth}}
\newcommand{\rank}{\operatorname{rank}}
\newcommand{\holm}{p_{\mathrm{Holm}}}

\title{Long-Horizon Predictive Subspaces in Mouse Secondary Motor Cortex Track Time in the Session, Not the Block Prior}

\author{%
  Gunnar Andra-Thomas\\
  Georgia Institute of Technology
}

\begin{document}

\maketitle

\begin{abstract}
Does the horizon of a predictive objective determine what a neural population subspace encodes? I fit lagged
reduced-rank ridge regression to Neuropixels recordings from mouse secondary motor cortex (MOs; 16 sessions, 13 mice,
International Brain Laboratory Brain-Wide Map) at lags from 0.5 to 30\,s, and asked whether long-lag subspaces carry
the Bayes-optimal block prior while short-lag subspaces carry the current stimulus. The plan and three primary tests
were pre-registered before any real-data fit. The main hypothesis was not supported. The prior showed no sign of
rising with lag (95\% upper bound on the mean change $+0.002$ in corrected $R^2$) and if anything fell; the only weak
evidence for it was in the 0.5\,s subspace. Current-stimulus information fell with lag, and 0.5\,s and 30\,s subspaces
differed by more than their split-half floors, but both effects are modest. Every five-dimensional summary was
dominated by a slow, near-linear mode that tracks time in the session: elapsed-time $R^2$ was 0.84 for the 30\,s
subspace, 0.82 at 0.5\,s, 0.80 for PCA and 0.64 for covariance-matched random subspaces. The pattern generalized to
neighbouring areas recorded on the same probes (15 sessions, 12 mice). Exploratory pre-registered follow-ups found no evidence that unit turnover, spike sorting, rigid probe motion or
spike-amplitude loss accounts for the slow mode (beyond an inconclusive five-session amplitude-stable subset), and
its fluctuations around the trend were neither shared across areas nor related to pupil, engagement, satiety,
reaction time or accuracy. Its origin remains open. In these recordings, long horizons recover session-scale drift
only slightly more than other low-dimensional summaries do, and such subspaces need elapsed-time diagnostics and
covariance-matched baselines before being read as task-related.
\end{abstract}

% =====================================================================================================
\section{Introduction}
\label{sec:intro}

Neural population activity mixes processes on many timescales, from sub-second sensorimotor transients to internal
states that drift over minutes. A natural way to separate them is by predictability: a direction in population space
that still predicts activity 30\,s later must carry a slow process, whereas one that predicts only 0.5\,s ahead can be
fast. This project asked whether the \emph{prediction horizon} of a linear predictive objective selects what the
resulting population subspace encodes. If it does, horizon-specific subspaces would offer an unsupervised route to
latent cognitive variables, such as beliefs about hidden task state, that are otherwise found only by supervised
decoding; if it does not, it matters to know what long-horizon objectives, such as time-lagged and slowness-based
methods \citep{perezhernandez2013,clark2019}, recover instead.

The International Brain Laboratory (IBL) visual decision task is well suited to the question
\citep{ibl2021task,ibl2025bwm}. On each trial a grating appears on the left or right, and the probability that it
appears on the left switches between 0.2 and 0.8 in uncued blocks lasting tens of trials. The Bayes-optimal block
prior, an ideal observer's running estimate of the current block given the stimuli seen so far, is a normative,
slowly varying latent variable that \citet{findling2025} decoded from many brain regions. Stimulus, choice and
movement, by contrast, vary within a trial.

The proposal stated two hypotheses. \emph{Geometry}: subspaces optimized to predict activity 30\,s ahead differ from
those optimized for 0.5\,s by more than the split-half noise floor of either. \emph{Content}: decoding of the prior
rises with the lag, while decoding of the current stimulus side falls. I tested both in secondary motor cortex (MOs),
the only candidate region with enough sessions, under a pre-registered plan (Section~\ref{sec:prereg}).

\paragraph{Summary of findings.} The prior is present in the full MOs population, replicating \citet{findling2025}. It
did not rise with the prediction horizon; if anything it fell, and the only weak evidence for it in a predictive
subspace was at 0.5\,s. The other two pre-registered predictions held, but weakly: long-lag subspaces carry less
current-stimulus information (because they lose it, not because short-lag subspaces specialize in it), and the 0.5\,s
and 30\,s subspaces differ by more than their split-half floors (mostly because short-lag subspaces are reproducible
and long-lag ones are not). Every low-dimensional summary, and the long-lag subspaces most of all, is dominated by one
slow mode that tracks time in the session. The same pattern appeared in neighbouring areas on the same probes.
Exploratory follow-ups found no evidence that the recording artifacts most likely to produce the mode account for it,
beyond an inconclusive five-session subset of amplitude-stable units, and no behavioural covariate explained its
fluctuations. The proposal's risk table listed this outcome (``the slow
subspace is drift: it encodes elapsed time, and prior decoding does not survive the null''), with the response
``report it as a negative result on long-lag objectives''; this report follows that row.

% =====================================================================================================
\section{Related work}
\label{sec:related}

\paragraph{Predictive and dynamical dimensionality reduction.} Reduced-rank regression (RRR; \citealp{izenman1975}),
here in its ridge-regularized form \citep{mukherjee2011}, finds low-dimensional linear maps between two sets of
variables; it defines communication subspaces between cortical areas \citep{semedo2019} and has been used to predict
later from earlier activity \citep{boucher2023}. Time-lagged independent component analysis recovers the slowest
processes of a dynamical system \citep{perezhernandez2013}, and Dynamical Components Analysis selects projections that
maximize past-future predictive information \citep{clark2019}. The canonical analysis of \citet{boxtiao1977} orders
components of a time series from least to most predictable; in a nonstationary series the most predictable ones are
the slowly wandering ones. For variance-based methods, random walks concentrate their variance in a few smooth,
time-ordered principal components \citep{antognini2018}, and PCA can produce oscillation-like components from smooth
data \citep{shinn2023}. The covariance-matched random-subspace baseline used here is that of \citet{elsayed2016}; see
\citet{elsayed2017} for the general case for such chance baselines.

\paragraph{The block prior in IBL data.} \citet{findling2025} decoded the Bayes-optimal prior from pre-stimulus
activity across the Brain-Wide Map \citep{ibl2025bwm}, using pseudo-sessions (block sequences resampled from the task's
generative process) to control for slow drift. That null follows \citet{harris2020}, who showed that
cross-validation, circular shifts and detrending can all give high false-positive rates when slowly drifting signals
are correlated, whereas tests built from the experimenter's own randomization remain valid. In the same dataset,
\citet{shi2025} mapped intrinsic timescales of single neurons, but not of population subspaces.

\paragraph{Slow drift, arousal and movement.} \citet{cowley2020} described a one-dimensional ``slow drift'' of
macaque V4 and prefrontal activity that covaried with impulsivity-like behaviour and pupil size; their evidence for
a genuine state rested on waveform-stability controls, agreement across separately recorded areas, and behavioural
covariation. In mice, infraslow population activity is coherent with pupil-linked arousal \citep{okun2019}, thirst
and satiety reorganize brain-wide activity within a session \citep{allen2019}, and pre-stimulus activity predicts
engagement \citep{steinmetz2019}. Engagement fluctuates within IBL sessions, as captured by GLM-HMM strategy states
(hidden Markov models whose states are separate generalized linear models of choice; \citealp{ashwood2022}). The
highest-variance dimensions of mouse cortical activity are dominated by spontaneous and uninstructed movements
\citep{stringer2019,musall2019}, and behavioural-state changes can masquerade as time-dependent drift \citep{sadeh2022}.

\paragraph{Recording nonstationarity.} Neuropixels probes \citep{jun2017} move relative to the brain, and IBL's spike
sorting corrects for this motion \citep{iblsort2022}; residual motion and amplitude change can make units fade in and out \citep{iblsort2022,iblrepro2025}, putting a slow, time-locked component into any population analysis. I found no prior work
reporting that the leading component of a long-horizon predictive subspace tracks session time.

% =====================================================================================================
\section{Data and methods}
\label{sec:methods}

\subsection{Data and inclusion}
\label{sec:data}

I used the public IBL Brain-Wide Map (BWM) release \citep{ibl2025bwm} (aggregate tag
\texttt{2024\_Q2\_IBL\_et\_al\_BWM}, spike-sorting revision 2024-05-06), keeping units that pass all of IBL's
single-unit quality metrics \citep{iblrepro2025} and trials that pass the BWM trial mask (reaction time between 0.08
and 2\,s, no missing task events). Regions follow the Allen-atlas ``Beryl'' parcellation used by the BWM. A census
rule fixed before analysis made a candidate region (MOp, MOs, ACAd, VISp) eligible if it had at least 15 sessions from
at least 8 mice meeting inclusion, and picked the eligible region with the largest published prior effect
\citep{findling2025}. MOs was the only eligible region (16 sessions, 13 mice; MOp 14 sessions from 12 mice, VISp 10/8,
ACAd 4/4); MOp's published effect was slightly larger (0.061 versus 0.052). Inclusion required at least 20 units
firing at least 0.1\,Hz over the task epoch and at least 400 masked trials. Sessions had 22 to 128 units (median 48.5;
842 in total) and 413 to 700 trials (median 534); the median task epoch lasted about 66 minutes. Spike counts were
binned in 0.5\,s bins over the task epoch, square-root transformed and z-scored per unit. Lags were
$\Delta \in \{0.5, 2, 8, 30\}$\,s.

\subsection{Lagged reduced-rank ridge regression}
\label{sec:rrr}

Let $x_t \in \RR^p$ be the z-scored population vector in bin $t$, and let the rows of $X$ and $Y$ hold $x_t$ and
$x_{t+\Delta}$ for all valid pairs, each column-centred. For each lag the subspace comes from
\begin{equation}
  \hat B_\Delta = \arg\min_{\rank(B)\le d} \; \lVert Y - XB \rVert_F^2 + \lambda \lVert B \rVert_F^2 .
  \label{eq:rrr}
\end{equation}
This has a closed form \citep{mukherjee2011}. Let $\hat B_R = (X^\top X + \lambda I)^{-1}X^\top Y$ be the full-rank
ridge solution, and let $V_d \in \RR^{p\times d}$ hold the top $d$ eigenvectors of
\begin{equation}
  G = \hat B_R^\top X^\top X \hat B_R + \lambda\, \hat B_R^\top \hat B_R ,
  \label{eq:gram}
\end{equation}
the top right singular vectors of the stacked matrix $[X\hat B_R;\, \sqrt{\lambda}\,\hat B_R]$. Then
\begin{equation}
  \hat B_\Delta = \hat B_R V_d V_d^\top, \qquad
  U_\Delta = \orth(\hat B_R V_d), \qquad
  z^{(\Delta)}_t = U_\Delta^\top x_t .
  \label{eq:subspaces}
\end{equation}
The \emph{readout basis} $U_\Delta$ spans the directions of current activity that predict the future, and the
\emph{loading basis} $V_d$ spans the predicted activity; primary analyses use the readout basis.

The penalty is $\lambda = \alpha\,\tr(X^\top X)/p$, with $\alpha$ on 13 log-spaced values from $10^{-3}$ to $10^{3}$.
Because the penalty itself changes which subspace is found, I chose \emph{one} $\alpha$ per session (the value
maximizing mean held-out rank-$d$ $R^2$ across the four lags, from five contiguous folds with a 10-bin embargo and
fold-wise z-scoring), so the penalty cannot make lags differ. The rank was $d=5$, with $d=3$ and $d=10$ as
pre-registered sensitivity analyses ($d=10$ only in sessions with more than 40 units: 12 sessions, 10 mice).
Baselines at the same $d$ were PCA followed by ridge regression from the scores to future activity, and a
covariance-matched random subspace $\orth(E\Lambda^{1/2}G)$, where $E\Lambda E^\top$ is the data covariance and $G$ a
$p\times d$ standard Gaussian matrix \citep{elsayed2016}. Full-rank ridge is the unconstrained reference.

\subsection{Geometry}
\label{sec:geometry}

Each session was split into halves separated by a 30\,s gap, and every lag was refit on each half with the
whole-session z-scoring and the session $\alpha$. For orthonormal bases $Q_A$ and $Q_B$, the principal angles
$\theta_1\le\dots\le\theta_d$ satisfy $\cos\theta_k = \sigma_k(Q_A^\top Q_B)$, with $\sigma_k$ the $k$-th singular
value; the summary is the mean angle. The \emph{cross-lag} angle compares the 0.5\,s fit on one half with the 30\,s fit
on the other (averaged over both assignments); the \emph{same-lag floors} compare one lag's fits on the two halves.
Chance is the mean angle between two independent covariance-matched random subspaces. The \emph{nesting index}
$\nu = \lVert Q_A^\top b \rVert^2$ is the fraction of a subspace's leading direction $b$ inside another subspace
$\operatorname{span}(Q_A)$; it is 1 for full containment and about $d/p$ for a random direction.

\subsection{Content decoding and the pseudo-session null}
\label{sec:decoding}

Trial-aligned spike counts in fixed 0.5\,s windows were transformed with the session z-scoring and projected onto
each basis. The prior and previous stimulus side were decoded from the pre-stimulus window (600 to 100\,ms before
stimulus onset, as in \citealp{findling2025}); the current stimulus side from the 0.5\,s after onset, on
nonzero-contrast trials; choice from the 0.5\,s before first movement; feedback from the 0.5\,s after feedback.
Binary targets used $L_2$ logistic regression ($C=1$) scored by mean per-fold area under the ROC curve (AUC);
continuous targets used ridge regression ($\alpha=1$) scored by pooled out-of-fold $R^2$, with five contiguous folds
and settings identical across lags, real data and null. Wheel speed and elapsed time were decoded from the continuous
bins, with 20 contiguous chunks assigned round-robin to five folds and a 10-bin embargo. Because the subspaces are fit
on the full session and each test chunk sits between training chunks, elapsed-time $R^2$ measures interpolation: the
share of a subspace's activity that varies smoothly with session time, not the ability to extrapolate it.

For the experimenter-generated variables (prior, previous and current side), each score was compared with the same
decoder applied to 200 pseudo-sessions, stimulus sequences resampled from the task's generative process with the
prior recomputed for each \citep{harris2020,findling2025}. The corrected score is the real score minus the null
median. The same pseudo-sequences were used for every subspace and lag, so a lag contrast is corrected by the null of
the \emph{same paired difference}. With contiguous folds and a slowly varying target, raw prior $R^2$ from the
subspaces is negative (session means $-0.126$ at 0.5\,s to $-0.287$ at 30\,s, PCA $-0.108$); because pseudo-sessions
pass through the identical decoder and folds, the corrected score still measures association beyond chance block
sequences, but on that decoder's scale rather than as a variance fraction. Elapsed time is not randomized, so it has no
null; a session's \emph{drift flag} (the proposal's rule) is set when the 30\,s subspace decodes elapsed time better
than both PCA and the random subspace while its prior decoding is not significant. A separate \emph{prior gate}
decoded the prior from the full population with Findling-style Lasso and interleaved nested cross-validation against
200 pseudo-sessions, combining sessions by Fisher's method; its scale therefore differs from the subspace decoders'.
Further details are in Appendix~\ref{app:impl}.

\subsection{Statistics, pre-registration and amendments}
\label{sec:prereg}

Sessions were averaged within each mouse, and conditions were compared with one-sided Wilcoxon signed-rank tests
across the 13 mice (smallest attainable exact $p$ $2^{-13}\approx1.2\times10^{-4}$). Confidence intervals are percentile
bootstraps over mice, which under-cover with $n=13$. Unless stated, summaries are means over mice. The three primary
tests, Holm-corrected as a family \citep{holm1979}, were: \textbf{P1 (geometry)}, the cross-lag angle exceeds the
same-lag floor at 0.5\,s \emph{and} at 30\,s (an intersection-union test, $p$ = the larger one-sided $p$);
\textbf{P2 (prior)}, the corrected prior $R^2$ rises from 0.5\,s to 30\,s; \textbf{P3 (current side)}, the corrected
current-side AUC falls from 0.5\,s to 30\,s. ``Pre-registered'' here means written in a timestamped commit to the
author's private repository before the relevant data were analysed, not a public registry. The amendment file,
committed before any real-data fit, lists nine deviations from the proposal and eleven additions
(Appendix~\ref{app:prereg}), including a pre-registered expectation that a shared slow mode would make P2 null and a
trial-locked residual rerun (event responses removed by a linear GLM with raised-cosine kernels, \citealp{pillow2008})
as a sensitivity analysis.

\paragraph{A post-registration bug fix.} The z-scoring floor on the standard deviation ($10^{-3}$) let a unit silent
across a training fold but active in its test fold dominate the pooled held-out variance, corrupting cross-validated
$R^2$ and the chosen $\alpha$. The floor was raised to 0.218, the standard deviation of the square-root count of a
unit at the 0.1\,Hz inclusion rate (amendment D10). It changed $\alpha$ in 2 sessions (4 in the residual run) and no
whole-session z-score. Raw $d=5$ primary and leave-one-out $p$-values were unchanged; these moved: residual P1
0.048 to 0.013, residual P3 0.012 to 0.0051, residual worst leave-one-mouse-out P1 0.092 (not significant) to 0.027
and P3 0.024 to 0.010, $d=3$ P1 0.049 to 0.032; the P2 values stayed far from significance. An earlier estimator
result (reduced rank beating full-rank ridge at 30\,s) was an artifact of the bug. All numbers here are post-fix.

The proposal froze analyses on 20 November (``bug fixes only from here''). Leave-one-mouse-out refits are post hoc.
The four follow-up pilots (Sections~\ref{sec:region2} and~\ref{sec:slowmode}) are exploratory, uncorrected, and do not
change the main analysis; the first was registered on 17 November, the other three on 21 and 22 November, after the
freeze.

% =====================================================================================================
\section{Results}
\label{sec:results}

\subsection{The prior is in MOs, and lagged RRR finds predictable structure}
\label{sec:gate}

The prior gate passed: the median corrected prior $R^2$ from the full population was 0.064, 5 of 16 sessions were
individually significant, and the Fisher-combined $p$ was $1.2\times10^{-5}$. Lagged RRR predicted future activity
better than PCA at every lag, by 0.011 at 0.5\,s up to 0.022 at 30\,s (all $p\le0.0053$, uncorrected secondary
tests), and at 30\,s five dimensions recovered 96\% of the full-rank held-out $R^2$ (Appendix
Table~\ref{tab:estimator}): structure predictable 30\,s ahead is very low-dimensional. After removing trial-locked
responses, RRR's margin over PCA shrank to 0.005 to 0.008 and its confidence interval included zero at every lag, so
much of the predictive objective's advantage over variance is task structure that recurs on every trial.

\subsection{Primary tests}
\label{sec:primary}

\begin{table}[t]
  \caption{Pre-registered primary tests ($d=5$, 13 mice; one-sided Wilcoxon across mice, Holm over three tests).
  Effect: P1, cross-lag angle minus the 0.5\,s and 30\,s floors (degrees); P2, change in corrected prior $R^2$ from
  0.5\,s to 30\,s; P3, drop in corrected current-side AUC. Brackets: 95\% bootstrap CI over mice. Mice agreeing: mice
  whose effect has the predicted sign.}
  \label{tab:primary}
  \centering\small
  \begin{tabular}{llccc}
    \toprule
    Test & Data & Effect [95\% CI] & $\holm$ & Mice agreeing \\
    \midrule
    \multirow{2}{*}{P1 geometry} & raw      & $+12.7$ [10.2, 15.1] and $+3.6$ [2.2, 4.9] & 0.0026 & 11 / 13 \\
                                 & residual & $+11.4$ [8.9, 13.9] and $+2.1$ [0.7, 3.3]  & 0.013  & 10 / 13 \\
    \multirow{2}{*}{P2 prior}    & raw      & $-0.099$ [$-0.218$, $+0.002$] & 0.89  & 4 / 13 \\
                                 & residual & $-0.093$ [$-0.154$, $-0.038$] & 0.996 & 3 / 13 \\
    \multirow{2}{*}{P3 current side} & raw      & $+0.061$ [0.028, 0.094] & 0.0061 & 10 / 13 \\
                                     & residual & $+0.063$ [0.032, 0.097] & 0.0051 & 11 / 13 \\
    \bottomrule
  \end{tabular}
\end{table}

\paragraph{Geometry (P1): significant but thin.} The mean principal angle between halves was $58.9^\circ$ for
0.5\,s versus 0.5\,s, $68.0^\circ$ for 30\,s versus 30\,s and $71.6^\circ$ across lags, against a covariance-matched
chance level of $67.9^\circ$ (Table~\ref{tab:primary}, Figure~\ref{fig:geometry}a). Most of the effect comes from the
short-lag side: the 30\,s floor is itself at chance, so ``cross exceeds the 30\,s floor'' compares against a noise-level
baseline, by $3.6^\circ$. P1 held with per-lag penalties, the sum of squared sines and the first angle (uncorrected
$p$ 0.0012, 0.0006, 0.013) but fails in the loading basis: there the cross-lag angle ($67.6^\circ$) is $1.0^\circ$
\emph{below} the 30\,s floor ($p=0.96$; Figure~\ref{fig:geometry}c), so the result is a property of the readout basis.
There is also a tension: the cross-lag angle is above chance even though the full-session subspaces are nested (the
0.5\,s subspace contains 92\% of the 30\,s leading direction; Section~\ref{sec:slow}). Nesting concerns one direction
of full-session fits, whereas the angles average all five dimensions of half-session fits; and chance is not a neutral
null here, because with the small penalties selected ($\alpha\le0.32$ in every session, one at the grid edge of 0.001)
the readout basis $\orth(\hat B_R V_d)$ is pushed toward low-variance input directions, while covariance-matched random
subspaces are biased toward high-variance ones. ``Above chance'' should not be read as ``more different than random''.

\begin{figure}[t]
  \centering
  \maybefig{rfig_geometry.pdf}{width=\linewidth}{1.9in}
  \caption{Subspace geometry ($d=5$; mouse means, as in the tests). \textbf{a}, readout basis: per-mouse (grey) and
  mean (blue; values in parentheses) angles between fits on opposite halves; dashed, covariance-matched chance. The low grey trace is DY\_014
  (Section~\ref{sec:robust}). \textbf{b}, all five principal angles, readout basis. \textbf{c}, the same comparisons in
  the loading basis, where the cross-lag angle does not exceed the 30\,s floor.}
  \label{fig:geometry}
\end{figure}

\paragraph{Prior (P2): no rise with lag; weak evidence only at 0.5\,s.} There was no sign of a rise: the mean change
from 0.5\,s to 30\,s was $-0.099$ (median over mice $-0.041$), and the 95\% upper bound, $+0.002$ on this decoder's
corrected-$R^2$ scale, leaves little room for one. After residualization the change was significantly negative
($-0.093$ [$-0.154$, $-0.038$]); this reverse-direction result is post hoc and uncorrected. Per subspace
(Figure~\ref{fig:content}a), the corrected prior $R^2$ was 0.054 [0.006, 0.105] at 0.5\,s, then $-0.027$, $-0.057$ and
$-0.046$ at 2, 8 and 30\,s, against 0.013 for PCA and $-0.030$ for the random subspace. The 0.5\,s value excludes zero
across mice (0.053 [0.005, 0.104] after residualization), but the across-session Fisher combination is not significant
($p=0.15$; 0.050 residual; 0.12 to 0.27 at the longer lags), so this is weak evidence. After residualization PCA also
carries some prior (0.041 [0.010, 0.072]). The prior is thus weakly present at short lags and in variance-ranked
directions, and not detected in the long-lag subspaces; the 30\,s-only complement carries 0.005.

\paragraph{Current stimulus side (P3): long lags lose it; short lags are not specialized.} Corrected current-side AUC
fell from 0.218 at 0.5\,s to 0.144 to 0.161 at the longer lags (Figure~\ref{fig:content}b). The drop persists after
controlling for block information (pre-registered addition A5): the AUC gain over a prior-only model falls from
0.074 at 0.5\,s to 0.028 to 0.044. But the 0.5\,s subspace equals a covariance-matched random subspace (0.218) and is
below PCA (0.240): long lags lose stimulus information, as evoked responses lasting under a second would predict,
rather than short lags gaining it. P3 held at $d=10$ but not at $d=3$ ($\holm=0.11$).

\begin{figure}[t]
  \centering
  \maybefig{rfig_content.pdf}{width=\linewidth}{1.9in}
  \caption{What the subspaces encode, by RRR lag, with PCA and random baselines (orange). Grey: individual mice;
  colour: mean over mice with 95\% bootstrap CI. \textbf{a}, \textbf{b}: null-corrected (real minus pseudo-session
  median). \textbf{c}: held-out $R^2$ (interpolation; Section~\ref{sec:decoding}). All variables, including residual
  versions, are tabulated in Appendix Table~\ref{tab:content}.}
  \label{fig:content}
\end{figure}

\subsection{One shared slow mode that tracks time in the session}
\label{sec:slow}

Elapsed time was strongly decodable from every five-dimensional summary (Figure~\ref{fig:content}c): held-out $R^2$
0.84 for the 30\,s subspace and 0.82 for the 0.5\,s subspace, against 0.80 for PCA and 0.64 for covariance-matched
random subspaces (0.73, 0.72, 0.70 and 0.56 after residualization). Long lags add only a little over PCA, though the
increase with lag is consistent (per-session Spearman correlation with lag, mean 0.88, $p=1.2\times10^{-4}$,
secondary). The drift flag fired in 14 of 16 sessions (12 of 16 after residualization). The time signal is shared
rather than split off: the 0.5\,s subspace contains 92\% [90, 94] of the 30\,s subspace's leading direction (reverse,
76\%), and the directions unique to the 30\,s subspace carry little of it (elapsed-time $R^2$ 0.12 raw, 0.03 residual).
This is why P2 is null, as the pre-registered expectation note anticipated: the slow mode fills the 30\,s dimensions.

I call the leading direction of the full-session 30\,s readout subspace the \emph{slow mode}. Averaged in 30\,s
windows and sign-aligned to increase with time, it is close to a straight line in most sessions (median correlation
with time 0.96; above 0.9 in 12 of 16). Not always smoothly, though: of six example sessions chosen by a fixed rule
(Appendix Figure~\ref{fig:examples}), three are near-linear ramps, one rises in several steps, and two are flat with a
single abrupt step (near 49 and 43 minutes), which look more like discrete events than a gradual process.

\subsection{Robustness}
\label{sec:robust}

The primary results were stable (Appendix Table~\ref{tab:robust}). The prior test was never significant (smallest
Holm $p$ 0.21, at $d=3$); the current-side test failed only at $d=3$; the geometry test survived every variant,
including the worst leave-one-mouse-out refit ($\holm=0.0051$ raw, 0.027 residual). One mouse, DY\_014, is a low
outlier in Figure~\ref{fig:geometry}a (an exploratory check found its readouts dominated by one cluster firing at
about 141\,Hz, likely multi-unit); dropping it (post hoc) changes nothing material.

\subsection{Generalization to neighbouring areas on the same probes}
\label{sec:region2}

The MOs probes also pass through other areas. In a pre-registered extension (A15), for each session I took the non-MOs
area with the most good units on the same probes and reran the unchanged pipeline, with criteria fixed before these
data were binned. This gave 15 sessions from 12 mice and nine areas: frontal cortex (ACAd, ILA, PL, ORBl, ORBvl, AIv,
MOp) in 10 sessions and lateral septum or striatum (LSr, CP) in 5. All four pre-registered criteria held
(Table~\ref{tab:region2}), the rule's ``replicates'' reading; I describe it as generalization to neighbouring areas
recorded on the same probes. Two of the criteria are lenient: ``P2 not significant'' is a null, and ``elapsed-time
$R^2\ge0.5$'' is passed even by covariance-matched random subspaces (0.56). The prior is barely present in these
populations (median corrected $R^2$ 0.001; the gate, which ran after the fits and statistics, was significant overall),
so this area set supports ``no rise with lag'' but cannot speak to where the prior lives.

\begin{table}[t]
  \caption{Pre-registered generalization test in neighbouring areas on the same probes (A15), against MOs.
  Criteria: P1 and P3 significant, P2 not significant, 30\,s elapsed-time $R^2\ge0.5$.}
  \label{tab:region2}
  \centering\small
  \begin{tabular}{lcc}
    \toprule
     & MOs & Neighbouring areas \\
    \midrule
    Sessions / mice                                              & 16 / 13 & 15 / 12 \\
    Prior gate: median corrected $R^2$ (Fisher $p$)              & 0.064 ($1.2\times10^{-5}$) & 0.001 (0.002) \\
    P1 geometry ($\holm$); cross minus floors, 0.5\,s / 30\,s     & 0.0026; $12.7^\circ$ / $3.6^\circ$ & 0.0037; $15.0^\circ$ / $3.7^\circ$ \\
    P2 prior rises with lag ($\holm$)                            & 0.89 & 0.52 \\
    P3 current side falls with lag ($\holm$)                     & 0.0061 & 0.012 \\
    Elapsed-time $R^2$: 30\,s / 0.5\,s / PCA / random            & 0.84 / 0.82 / 0.80 / 0.64 & 0.82 / 0.78 / 0.71 / 0.56 \\
    Drift flags; nesting of 30\,s direction in 0.5\,s subspace   & 14/16; 0.92 & 14/15; 0.92 \\
    \bottomrule
  \end{tabular}
\end{table}

% =====================================================================================================
\section{What is the slow mode?}
\label{sec:slowmode}

A slow mode that tracks session time could be a recording artifact or a slow change in brain state. Unit turnover,
drifting spike sorting, rigid probe motion and spike-detection loss each predict a specific signature, as do arousal,
engagement and satiety. Four exploratory pilots tested these in turn, each pre-registered with a reading rule (summary
in Appendix Table~\ref{tab:pilots}). Their registered readings were ``mixed or undetermined'' (first), ``sorting
artifact'' (second; a label whose implied mechanism the third pilot's registered test then contradicted), ``amplitude
cannot account for the mode'' (third) and ``replicates'' (fourth). Figure~\ref{fig:controls} shows elapsed-time
decoding after refitting on each modified data copy.

\paragraph{Removing unstable units; probe motion; pupil (A12).} In 15 of 16 sessions at least one unit falls below
0.1\,Hz in some fifth of the session (257 of 842 units; 68 units in 13 sessions are completely silent for a fifth).
Keeping only units at or above 0.1\,Hz in every fifth left 13 sessions with at least 20 units; refit, the 30\,s
subspace still decoded elapsed time at 0.77 (PCA 0.72, random 0.50; the same sessions originally 0.83), and the
primary results were unchanged. IBL's motion estimate for these insertions is small (median range over a session
7.9\,$\mu$m). The slow mode correlates as strongly with \emph{other} sessions' motion and pupil traces as with its own
(session-permutation corrected scores $-0.006$ and $-0.001$, $p=0.39$ and 0.42). Because the mode is nearly a straight
ramp, such tests cannot discriminate causes that do not differ in shape.

\paragraph{Sorting-free activity, probe shift and spike amplitudes (A13).} Binning every detected spike, whatever its
cluster, into 50\,$\mu$m depth channels reproduced the mode (median correlation with the unit-based mode 0.96;
elapsed-time $R^2$ 0.74 against 0.88 for the same sessions), so I found no evidence that cluster assignment produces
it. Multi-unit loadings did not follow the depth-profile slope ($p=0.25$), which is absence of evidence for a rigid
shift rather than evidence against it. Spike amplitudes are far from stable (about 59\% of units have a consistent,
mostly declining trend), and units whose amplitude shrinks tend slightly to be those losing rate along the mode
(Fisher $z$ 0.15, $p=0.029$). Only 5 sessions kept 20 or more units stable in both rate and amplitude; there,
elapsed-time $R^2$ fell to 0.49 against 0.73 for the same sessions originally (67\%), but PCA also fell (0.68 to 0.43),
so part of the drop reflects fewer units; the rule read this ``intermediate''. Engagement (per-trial engaged-state
probability from IBL's GLM-HMM fits, \citealp{iblengage}) and satiety (cumulative reward) were not tied to the mode
beyond a quadratic trend ($p=0.19$, 0.12). The rule's ``sorting artifact'' label came from its second branch
(amplitude-stable refit not ``survives'' and the amplitude-loading test significant), but the multi-unit refit
contradicts a sorting cause; the rule had not anticipated amplitude loss acting through spike \emph{detection}.

\paragraph{Can amplitude loss account for the mode? (A14).} Regressing a unit's rate on its amplitude over time cannot
fail as a test, because both are near-ramps. The pilot instead fit, per unit, a Poisson GLM of spike count per 30\,s
window on relative amplitude $\tilde a_u(t)$ with a quadratic in time as nuisance terms. Detection loss predicts a
positive coefficient; it was negative (mean $-1.54$, $p=0.99$ for positive): a unit fires about 16\% \emph{more} in
windows where its spikes are 10\% smaller, consistent with the known shrinkage of spikes during fast firing and bursts
\citep{quirk1999,harris2000}, and this held for isolated spikes (post hoc). Corrected copies rescaled counts by
\begin{equation}
  f_u(t) = \exp\!\big(-\beta\, \tilde a_u(t)\big), \qquad f_u(t)\in[0.5,\,2] .
  \label{eq:ampcorr}
\end{equation}
The registered conservative arm used the session-median fluctuation coupling; because that $\beta$ is negative (median
$-1.43$), it corrects in the opposite direction from detection loss and does not test it. The informative arms use
each unit's own coupling with no time terms (maximal) and, post hoc, the same on a heavily smoothed amplitude
trajectory: they kept elapsed-time decoding at 95\% and 93\% of the original (all 16 sessions) and about 85\% of the
slope along the original slow direction. Amplitude change was not concentrated early, as settling after insertion
would predict (index $-0.12$, $p=0.75$). Caveats: about a third of units hit the clip in some window; spike-sorter
amplitudes are template scalings, not voltages; and settling before the task began would be invisible.

\begin{figure}[t]
  \centering
  \maybefig{rfig_controls.pdf}{width=0.92\linewidth}{2.3in}
  \caption{Held-out elapsed-time $R^2$ of the 30\,s subspace after refitting the full pipeline on each modified data
  copy. Grey: mice; filled marker and bar: mean over mice and 95\% bootstrap CI (orange: post hoc variant); open
  circle: the original value for the same sessions (post hoc comparison). The registered conservative amplitude
correction is omitted because it corrects in the wrong direction ($\beta<0$); solid line: original 0.84; dashed: the
  pre-registered 75\% (``survives'') and 50\% (``removed'') thresholds, defined relative to 0.84. Right: copy versus
  same-session original, with sessions/mice.}
  \label{fig:controls}
\end{figure}

\paragraph{Shared across areas, or tied to behaviour? (A15).} Before detrending, the MOs and neighbouring-area slow
modes correlate strongly (median 0.94), as any two ramps do. After a quadratic trend is removed from each, their
fluctuations are not shared beyond the cross-session null (median detrended $r$ 0.07; $p=0.48$, 12 mice). All the
areas were on the same probe as the MOs units, so this argues against a probe-level artifact as much as against a
shared state, though with 12 mice and small residual fluctuations the test has limited power. Beyond the trend, the
MOs slow mode tracked neither log reaction time nor accuracy ($p=0.12$, 0.34; signs inconsistent across mice).

\paragraph{Where this leaves the slow mode.} What remains is a slow, mostly near-linear change in firing that tracks
session time in both area sets. I found no evidence that unit turnover, spike sorting, rigid probe motion or
amplitude loss accounts for it, apart from the inconclusive five-session amplitude-stable subset; nor is it identified
as a brain state, since its fluctuations are not shared across areas and track no available behavioural measure. A
neural origin (slow arousal or neuromodulatory drift, fatigue, elapsed-time coding) remains possible but not
established; the ramp shape limits correlational tests, the step-like sessions suggest discrete events of unknown
cause, and unmeasured recording factors remain possible.

% =====================================================================================================
\section{Discussion}
\label{sec:discussion}

\paragraph{What the results say about the question.} The horizon changes which directions are selected, but not as
hypothesized. At 0.5\,s the subspace is reproducible and carries within-trial variables about as well as a random
covariance-matched subspace; at 30\,s it is dominated by a slow mode that is shared with the short-lag subspace, not
reproducible across halves, and indexes time rather than the latent variable the mouse must track. The prior does not
move into long-horizon subspaces; it is weakly present at short lags and, after residualization, in PCA directions.
Whether it is carried as a small, event-updated modulation in low-variance directions cannot be decided here: the
prior effect is weak in MOs and almost absent in the neighbouring areas, and absence from a rank-5 subspace filled by a
stronger slow signal does not show that the prior is non-dynamical.

\paragraph{A methodological lesson, specific to these data.} In these recordings every low-dimensional summary is
dominated by session-scale drift (PCA 0.80, random 0.64), and long horizons only slightly more so (0.84 at 30\,s, 0.82 at 0.5\,s). Long-horizon
objectives inherit the slowest variance of a session \citep{boxtiao1977}; elapsed-time decoding next to PCA and
covariance-matched baselines, pseudo-session nulls and care with fold-wise normalization are needed before reading
such subspaces as cognition.

\paragraph{Limitations.} This is one dataset of acute recordings, in 13 mice, and the second area set shares probes,
sessions and mice with MOs. Ramp-shaped signals inflate correlations with anything that trends, elapsed-time decoding
interpolates, and the subspaces are fit on all of each session's data. Decoders are linear and the prior is read from
one pre-stimulus window, so a nonlinear or differently timed prior code would be missed. Percentile bootstrap intervals
with $n=13$ under-cover, P3 is sensitive to $d$, P1 is thin, and the pilots are exploratory and uncorrected. The
detrended tests (cross-area sharing, behaviour, pupil, motion) have limited power, and a 30\,s lag leaves only about
120 independent windows per session.

\paragraph{Future work.} Because the core pattern generalizes without pipeline changes, the next step is independent
sessions and regions, which needs new downloads (MOp has 14 sessions from 12 mice). Spontaneous periods outside the
task could show whether the slow mode is task-linked.

\paragraph{Conclusion.} Predicting population activity further ahead did not isolate the block prior, in MOs or in
neighbouring areas; it isolated a slow, time-tracking mode, slightly more strongly than other low-dimensional
summaries, whose origin none of the controls tested here resolved.

\section*{Code and data availability}
All recordings are from the public IBL Brain-Wide Map release \citep{ibl2025bwm} (aggregate tag
\texttt{2024\_Q2\_IBL\_et\_al\_BWM}, spike-sorting revision 2024-05-06), accessed with the IBL ONE interface. The
analysis code (library and numbered pipeline and pilot scripts, with unit tests and synthetic controls), the
pre-registration and amendments (\texttt{AMENDMENTS.md}) and the run-by-run analysis log (\texttt{ANALYSIS\_LOG.md})
are in the author's project repository and are available on request; result files are regenerated by the numbered
scripts.

% AI-use statement drafted for the author's review; edit to reflect your own contribution before submission.
% Author must confirm: which design decisions/amendments were yours vs suggested by Claude's adversarial review;
% what you wrote vs what Claude wrote.
\section*{AI use}
Claude (Anthropic), run as Claude Code, was used extensively in this project: to plan the analysis and adversarially
review the plan; to write the analysis code and tests; to run the analyses, locally and on the PACE ICE cluster; to
carry out the literature review; and to draft this report, the presentation slides and the figures, with further AI
agents reviewing the drafts. The author set the research question and wrote the proposal, made the design decisions
and amendments, and checked the results.

\bibliographystyle{plainnat}
\bibliography{references}

% =====================================================================================================
\clearpage
\appendix

\section{Pre-registration and amendments}
\label{app:prereg}

The amendment file was committed before any real-data fit. Table~\ref{tab:amend} summarizes every entry.
\emph{Pre-reg.} means written in a timestamped commit before the relevant data were analysed; \emph{post hoc} means
added after results were seen.

\begin{table}[H]
  \caption{Deviations from the proposal (D), additions (A), and later amendments.}
  \label{tab:amend}
  \centering\footnotesize
  \begin{tabularx}{\linewidth}{@{}l>{\raggedright\arraybackslash}X>{\raggedright\arraybackslash}X>{\raggedright\arraybackslash}p{2.3cm}@{}}
    \toprule
    \# & Change & Reason & Status \\
    \midrule
    D1 & Bin the task epoch only, not the whole session & Excludes passive and spontaneous periods & pre-reg. \\
    D2 & $\ge 400$ trials counted after the BWM trial mask & Unambiguous, stricter & pre-reg. \\
    D3 & Bayes-optimal prior without skipping no-go trials & Prior depends only on earlier stimuli & pre-reg. \\
    D4 & One cross-validated $\alpha$ per session, shared by all lags & Removes the $\lambda$ confound & pre-reg. \\
    D5 & Current side decoded on nonzero-contrast trials only & Zero-contrast side carries only block information & pre-reg. \\
    D6 & One-sided Wilcoxon tests & Directional hypotheses & pre-reg. \\
    D7 & Geometry as intersection-union test; Holm over 3 tests & Proposal required both floors & pre-reg. \\
    D8 & Halves use whole-session z-scoring and the session $\alpha$ & Same coordinates in both halves & pre-reg. \\
    D9 & Drop units below 0.1\,Hz & Guards z-scoring & pre-reg. \\
    A1 & Full-population prior gate (Lasso, pseudo-sessions) & Go/no-go before main analyses & pre-reg. \\
    A2 & Covariance-matched random baseline & Chance level for angles and estimator & pre-reg. \\
    A3--A4 & Nesting index; orthogonal-complement and single-dimension decoding & Distinguish nested from absent & pre-reg. \\
    A5 & Side variables as AUC gain over a prior-only model & Side variables are confounded with the block & pre-reg. \\
    A6 & Null shared across lags; paired-difference null; real minus null median (proposal: mean) & Valid lag contrasts; as Findling et al. & pre-reg. \\
    A7 & Bootstrap CIs; Fisher-combined session $p$; Spearman trend over lags & Secondary summaries & pre-reg. \\
    A8 & Purge, 10-bin embargo; 60-bin gap between halves & Leakage hygiene & pre-reg. \\
    A9 & Choice and feedback probe windows & Unspecified in proposal & pre-reg. \\
    A10 & Drift flag coded explicitly & Proposal's rule made explicit & pre-reg. \\
    A11 & Trial-locked residual rerun & Sensitivity & pre-reg. \\
    D10 & z-score SD floor raised from $10^{-3}$ to 0.218 & Bug (Section~\ref{sec:prereg}) & post hoc fix \\
    P1h & Leave-one-mouse-out and flagged-session robustness & Robustness only & post hoc \\
    A12 & Pilot: stable units, probe motion, pupil (17 Nov) & What is the slow mode? & pre-reg., exploratory \\
    A13 & Pilot: engagement, satiety, amplitude, multi-unit (21 Nov) & Ramp-independent evidence & pre-reg., exploratory \\
    A14 & Pilot: amplitude-rate coupling, corrected refits, timing (22 Nov) & Is amplitude loss sufficient? & pre-reg., exploratory; 2 post hoc checks \\
    A15 & Neighbouring areas, cross-area slow mode, behaviour (22 Nov) & Does it generalize? & pre-reg., exploratory; completed \\
    \bottomrule
  \end{tabularx}
\end{table}

\paragraph{Relation to the proposal.} The proposal promised three figures (geometry, the variables-by-lags decoding
matrix and the estimator check); they appear here as Figure~\ref{fig:geometry}, Figure~\ref{fig:content} with
Table~\ref{tab:content}, and Figure~\ref{fig:estimator}. Its risk table had a row ``the slow subspace is drift: it
encodes elapsed time, and prior decoding does not survive the null'', with the response ``report it as a negative
result on long-lag objectives''; this report follows it. A separate row, ``no dissociation across lags: Analysis 1 shows
the subspaces coincide'', with the response ``report the shared slow mode and characterize what it encodes'', was not
triggered (P1 was significant), although Section~\ref{sec:slowmode} does that characterization. Of the proposal's
ordered extensions, a second region (here, neighbouring areas in the same sessions), the GLM residual rerun and the
GLM-HMM engagement probe were carried out. The main analysis, including the rerun after the bug fix, was completed on
16 November 2026, before the 20 November freeze; the pilots do not alter it.

\section{Implementation details}
\label{app:impl}

\paragraph{Units, trials, regions.} Units are those passing all of IBL's single-unit quality metrics (BWM ``good''
units, label 1; \citealp{iblrepro2025}), spike-sorted with IBL's pipeline \citep{iblsort2022}. The BWM trial mask
excludes trials with reaction time outside 0.08 to 2\,s or with a missing stimulus, choice, feedback, block or
first-movement time. Regions use the Beryl mapping of the Allen atlas distributed with the BWM \citep{ibl2025bwm}.

\paragraph{Pseudo-sessions.} Block sequences follow the IBL generator: the first 90 trials have $P(\text{left})=0.5$;
thereafter blocks alternate between 0.2 and 0.8, with lengths drawn from an exponential with mean 60 and truncated to
integers from 20 to 99. Stimulus sides are drawn from each block probability, and the Bayes-optimal prior is
recomputed from the pseudo sides.

\paragraph{Trial-locked residuals.} For each event (stimulus onset, first movement, feedback), 8 raised-cosine basis
functions tile $-1$ to 2\,s around the event \citep{pillow2008}; summed responses form a design matrix fit to the
square-root counts by ridge regression (penalty 1). The design is not side- or block-specific, so only responses
identical on every trial are removed. The same coefficients are applied to the trial windows.

\paragraph{Penalty selection.} The chosen $\alpha$ was 0.1 in 9 sessions, 0.32 in 3, 0.032 in 2, 0.01 in 1 and 0.001
(the grid edge) in 1.

\paragraph{Unit stability.} Fifths are five equal parts of each session's valid task-epoch bins. Counts in
Section~\ref{sec:slowmode} were recomputed from the binned data with these definitions.

\section{Additional results}
\label{app:extra}

\begin{table}[H]
  \caption{Decoding from each subspace (means over 13 mice; raw data, then trial-locked residuals). AUC and side
  variables marked ``corr.'' are real minus pseudo-session median; the gain row is AUC gained over a prior-only model.}
  \label{tab:content}
  \centering\footnotesize
  \begin{tabular}{lcccccc}
    \toprule
    Variable & RRR 0.5\,s & RRR 2\,s & RRR 8\,s & RRR 30\,s & PCA & Random \\
    \midrule
    \multicolumn{7}{l}{\emph{Raw}} \\
    Current side (AUC, corr.) & 0.218 & 0.161 & 0.144 & 0.158 & 0.240 & 0.218 \\
    Current side, gain over prior (AUC) & 0.074 & 0.044 & 0.028 & 0.039 & -- & -- \\
    Previous side (AUC, corr.) & 0.045 & 0.050 & 0.044 & 0.037 & 0.027 & 0.007 \\
    Block prior ($R^2$, corr.) & 0.054 & $-$0.027 & $-$0.057 & $-$0.046 & 0.013 & $-$0.030 \\
    Elapsed time ($R^2$) & 0.821 & 0.834 & 0.839 & 0.841 & 0.800 & 0.641 \\
    Choice (AUC) & 0.615 & 0.604 & 0.589 & 0.590 & 0.638 & 0.616 \\
    Feedback (AUC) & 0.803 & 0.708 & 0.600 & 0.609 & 0.807 & 0.790 \\
    Wheel speed ($R^2$) & 0.339 & 0.191 & 0.024 & 0.043 & 0.406 & 0.272 \\
    \midrule
    \multicolumn{7}{l}{\emph{Trial-locked residuals}} \\
    Current side (AUC, corr.) & 0.224 & 0.161 & 0.169 & 0.162 & 0.271 & 0.245 \\
    Current side, gain over prior (AUC) & 0.076 & 0.043 & 0.044 & 0.041 & -- & -- \\
    Previous side (AUC, corr.) & 0.043 & 0.047 & 0.049 & 0.033 & 0.038 & 0.016 \\
    Block prior ($R^2$, corr.) & 0.053 & $-$0.017 & $-$0.034 & $-$0.038 & 0.041 & 0.004 \\
    Elapsed time ($R^2$) & 0.718 & 0.721 & 0.727 & 0.727 & 0.702 & 0.561 \\
    Choice (AUC) & 0.634 & 0.601 & 0.608 & 0.587 & 0.666 & 0.634 \\
    Feedback (AUC) & 0.767 & 0.685 & 0.633 & 0.654 & 0.802 & 0.761 \\
    Wheel speed ($R^2$) & 0.105 & 0.042 & 0.015 & 0.014 & 0.120 & 0.066 \\
    \bottomrule
  \end{tabular}
\end{table}

\begin{table}[H]
  \caption{Follow-up pilots (exploratory; uncorrected one-sided Wilcoxon across mice). ``$R^2$'': elapsed-time
  decoding by the 30\,s subspace after refitting on the modified data (original 0.84; same-session originals in
  Figure~\ref{fig:controls}).}
  \label{tab:pilots}
  \centering\footnotesize
  \begin{tabularx}{\linewidth}{@{}l>{\raggedright\arraybackslash}p{0.36\linewidth}>{\raggedright\arraybackslash}X@{}}
    \toprule
    Candidate & Test & Result \\
    \midrule
    Unit turnover (A12) & Refit on units at or above 0.1\,Hz in every fifth of the session & $R^2$ 0.77 (13 sessions, 12 mice); flag 12/13; reads ``survives'' \\
    Probe motion (A12) & Correlation with own vs.\ other sessions' IBL motion trace & Corrected $-0.006$ ($p=0.39$); detrended $+0.062$ ($p=0.14$) \\
    Pupil (A12) & Same, with pupil diameter (15 sessions) & Corrected $-0.001$ ($p=0.42$); detrended $-0.004$ ($p=0.55$) \\
    Spike sorting (A13) & Refit on sorting-free multi-unit activity in 50\,$\mu$m depth channels & $R^2$ 0.74 (15 sessions, 12 mice); match with unit-based mode 0.96 \\
    Rigid shift (A13) & Do multi-unit loadings follow the depth-profile slope? & Mean $|\rho|$ 0.19; Fisher $p=0.25$ \\
    Amplitude, loading (A13) & Do units whose amplitude shrinks carry the mode? & Mean Fisher $z$ 0.15, $p=0.029$ \\
    Amplitude, stable (A13) & Refit on rate- and amplitude-stable units & $R^2$ 0.49 (5 sessions, 5 mice); ``intermediate'' \\
    Engagement, satiety (A13) & Residual correlation beyond a quadratic trend, session-permutation null & $p=0.19$ and $p=0.12$; fluctuating-engagement subset $p=0.049$, signs inconsistent \\
    Detection coupling (A14) & Sign of amplitude-rate coupling beyond the trend & Negative: mean $-1.54$, $p=0.99$; degree-5 trend $-1.92$; isolated spikes (post hoc) $-1.32$, $p=0.98$ \\
    Amplitude correction (A14) & Refit after removing each unit's amplitude-linked rate & $R^2$ 0.80 / 0.78 (maximal / slow-maximal, post hoc), both ``survive''; the registered conservative arm (0.85) corrects in the wrong direction ($\beta<0$) and is not a test of detection loss \\
    Settling (A14) & Is amplitude change concentrated early? & Index $-0.12$ ($p=0.75$); if anything, later \\
    Shared across areas (A15) & Do MOs and neighbouring-area slow modes co-fluctuate beyond a quadratic trend? & Raw match 0.94; detrended corrected 0.006, $p=0.48$ (12 mice) \\
    Behaviour (A15) & Residual correlation with log reaction time and accuracy & $p=0.12$ and $p=0.34$; signs inconsistent \\
    \bottomrule
  \end{tabularx}
\end{table}

\begin{table}[H]
  \caption{Principal angles between session halves (degrees; means over 13 mice; readout basis).}
  \label{tab:angles}
  \centering\small
  \begin{tabular}{lcccc}
    \toprule
     & \multicolumn{2}{c}{Raw} & \multicolumn{2}{c}{Residual} \\
    Comparison & mean & first & mean & first \\
    \midrule
    0.5\,s vs 0.5\,s (same lag)         & 58.9 & 31.9 & 58.1 & 31.0 \\
    30\,s vs 30\,s (same lag)           & 68.0 & 48.2 & 67.5 & 43.7 \\
    0.5\,s vs 30\,s (cross lag)         & 71.6 & 52.3 & 69.5 & 46.8 \\
    Covariance-matched chance           & 67.9 & 47.0 & 68.8 & 48.6 \\
    \bottomrule
  \end{tabular}
\end{table}

\begin{table}[H]
  \caption{Robustness of the primary tests (Holm-corrected $p$). The $d$ variants and the residual rerun were
  pre-registered sensitivity analyses; leave-one-mouse-out and dropping DY\_014 are post hoc.}
  \label{tab:robust}
  \centering\small
  \begin{tabular}{lccc}
    \toprule
    Variant & P1 geometry & P2 prior & P3 current side \\
    \midrule
    Primary ($d=5$, raw)                         & 0.0026 & 0.89  & 0.0061 \\
    Trial-locked residuals                       & 0.013  & 0.996 & 0.0051 \\
    $d=3$                                        & 0.032  & 0.21  & 0.11   \\
    $d=10$ (12 sessions, 10 mice)                & 0.0029 & 0.88  & 0.0059 \\
    Leave one mouse out, worst case (raw)        & 0.0051 & 0.95  & 0.012  \\
    Leave one mouse out, worst case (residual)   & 0.027  & 0.998 & 0.010  \\
    Drop DY\_014 (raw)                           & 0.0051 & 0.87  & 0.0051 \\
    \bottomrule
  \end{tabular}
\end{table}

\begin{table}[H]
  \caption{Estimator check: held-out $R^2$ for predicting population activity $\Delta$ ahead ($d=5$; means over
  13 mice). RRR uses the session's shared penalty; the other methods use their best penalty.}
  \label{tab:estimator}
  \centering\small
  \begin{tabular}{lcccc}
    \toprule
    Method & $\Delta=0.5$\,s & $\Delta=2$\,s & $\Delta=8$\,s & $\Delta=30$\,s \\
    \midrule
    Full-rank ridge                 & 0.373 & 0.223 & 0.200 & 0.176 \\
    Lagged RRR                      & 0.246 & 0.192 & 0.183 & 0.172 \\
    PCA + ridge                     & 0.235 & 0.171 & 0.161 & 0.150 \\
    Covariance-matched random       & 0.162 & 0.133 & 0.121 & 0.118 \\
    \midrule
    RRR as fraction of full rank    & 0.65  & 0.82  & 0.88  & 0.96  \\
    \bottomrule
  \end{tabular}
\end{table}

\paragraph{Other secondary measures.} Corrected previous-side AUC was roughly flat across lags (0.037 to 0.050).
Current-side decoding from the 0.5\,s-only complement exceeded that from the 30\,s-only complement (0.259 versus 0.220;
$p=0.016$, uncorrected); prior decoding from the two complements did not differ ($p=0.39$). After residualization, the
nesting indices were 0.94 (0.5\,s containing 30\,s) and 0.89 (reverse), and RRR's margin over PCA was 0.005 to 0.008
with confidence intervals including zero at every lag. All three A14 corrected copies left the primary tests unchanged (geometry $\holm$ 0.0004 to 0.004, prior 0.63 to
0.73, current side 0.006 to 0.013). In the maximal and slow-maximal copies the drift flag fired in 14 of 16 sessions
and the corrected data kept 0.86 and 0.85 of the original slow-direction slope. The registered conservative copy
(13 of 16 flags, slope fraction 1.00) corrects in the wrong direction ($\beta<0$) and is not a test of detection
loss.

\clearpage
\section{Additional figures}
\label{app:figs}

\begin{figure}[H]
  \centering
  \maybefig{fig5_slowmode_examples.pdf}{width=0.85\linewidth}{2.1in}
  \caption{Slow-mode time courses (30\,s windows, z-scored, sign-aligned) in six sessions chosen by a fixed rule
  (first session of each of the first six mice in the frozen list), with straight-line fits and correlation $r$ with
  time. CSHL058, CSHL060 and KS023 are near-linear ramps; CSH\_ZAD\_024 rises in several steps (near 22--23, 55 and 80--82 minutes); CSHL049 and DY\_014 are flat with a single abrupt step (near 49 and 43 minutes).}
  \label{fig:examples}
\end{figure}

\begin{figure}[H]
  \centering
  \maybefig{fig3_estimator_d5.pdf}{width=0.8\linewidth}{2.6in}
  \caption{Estimator check: held-out $R^2$ of future activity against lag (mean $\pm$ s.e.m.\ over 13 mice).}
  \label{fig:estimator}
\end{figure}

\begin{figure}[H]
  \centering
  \maybefig{figS1_nesting_d5.pdf}{width=0.95\linewidth}{2in}
  \caption{\textbf{a}, nesting index of each subspace's leading direction in the other (dots: sessions; bar: mean;
  dashed: isotropic chance). \textbf{b}, null-corrected decoding from the orthogonal complements (mean and 95\% CI over
  mice).}
  \label{fig:nesting}
\end{figure}

\end{document}
